% GENERATED FILE. Edit canonical lessons, then run npm run generate:books.
% canonical-source-hash: fnv1a64:cf7ae45b0ec0ef89
% canonical-lessons: JA-C324-monku, JA-C324-kujou, JA-C324-fuhei, JA-C324-hihan-suru, JA-C324-iiwake

\chapter{Complaints}
\label{ch:ja-complaints}
\hlchaptermodality{324}

\begin{chapteropening}
\textbf{By the end of this chapter:} \emph{I can make a complaint and say what caused it.}

Complete the last lesson of chapter 324: I can make a complaint and say what caused it.
\end{chapteropening}

\section[\texorpdfstring{\ja{もんく} (monku)}{もんく (monku)}]{\ja{もんく} (monku) — a complaint; a grumble}

\label{lesson:JA-C324-monku}



\begin{quote}
Before the new one: say the Japanese for no thank you, then the Japanese for not particularly.
\end{quote}

\subsection*{You'll want to know: \ja{もんく}}
\textbf{\ja{もんく}} — \emph{monku} — \textquotedblleft{}a complaint; a grumble\textquotedblright{}.

\begin{quote}
\textbf{\ja{ねだんが} \ja{たかいから、きゃくが} \ja{もんくを} \ja{言}\ja{いました。}} — \emph{nedan ga takai kara, kyaku ga monku o iimashita} — \textquotedblleft{}The price was high, so the customer complained.\textquotedblright{}
\end{quote}

\begin{cousinweb}
\textbf{\ja{もんく}} (\emph{monku}) is Sino-Japanese (kango): two Chinese characters read with their on'yomi, readings borrowed from Chinese. They are \emph{mon}, \textquotedblleft{}sentence\textquotedblright{} (Middle Chinese \emph{mjun}) and \emph{ku}, \textquotedblleft{}phrase\textquotedblright{} (\emph{kjuH}) (KANJIDIC2; Middle Chinese in Baxter's transcription).
\end{cousinweb}

\begin{grammarlens}[title={What you've built}]
One more word for giving reasons and opinions.
\end{grammarlens}

\subsection*{Your turn}
\begin{itemize}
\raggedright
  \item \emph{Say it:} \emph{monku}
  \item \emph{Say it:} \emph{monku}, once more
  \item \emph{Say it:} \emph{nedan ga takai kara, kyaku ga monku o iimashita} — The price was high, so the customer complained.
  \item \emph{Recall:} say the Japanese for no thank you, then the Japanese for not particularly, then say \emph{monku} again
\end{itemize}

\begin{tcolorbox}[breakable,colback=teal!4,colframe=teal!35!black,title={Before you move on}]
What does \ja{もんく} mean? (\textquotedblleft{}A complaint; a grumble\textquotedblright{}.) Say it once more.
\end{tcolorbox}
\section[\texorpdfstring{\ja{くじょう} (kujō)}{くじょう (kujō)}]{\ja{くじょう} (kujō) — a formal complaint}

\label{lesson:JA-C324-kujou}



\begin{quote}
Before the new one: say the Japanese for not particularly, then the Japanese for a grumble.
\end{quote}

\subsection*{You'll want to know: \ja{くじょう}}
\textbf{\ja{くじょう}} — \emph{kujō} — \textquotedblleft{}a formal complaint\textquotedblright{}.

\begin{quote}
\textbf{\ja{となりの} \ja{おとが} \ja{やかましいので、くじょうを} \ja{言}\ja{いました。}} — \emph{tonari no oto ga yakamashii node, kujō o iimashita} — \textquotedblleft{}The noise next door was loud, so I made a complaint.\textquotedblright{}
\end{quote}

\begin{cousinweb}
\textbf{\ja{くじょう}} (\emph{kujō}) is Sino-Japanese (kango): two Chinese characters read with their on'yomi, readings borrowed from Chinese. They are \emph{ku}, \textquotedblleft{}suffering\textquotedblright{} (Middle Chinese \emph{khuX}) and \emph{jō}, \textquotedblleft{}circumstances\textquotedblright{} (\emph{dzjeng}) (KANJIDIC2; Middle Chinese in Baxter's transcription).
\end{cousinweb}

\begin{grammarlens}[title={What you've built}]
One more word for giving reasons and opinions.
\end{grammarlens}

\subsection*{Your turn}
\begin{itemize}
\raggedright
  \item \emph{Say it:} \emph{kujō}
  \item \emph{Say it:} \emph{kujō}, once more
  \item \emph{Say it:} \emph{tonari no oto ga yakamashii node, kujō o iimashita} — The noise next door was loud, so I made a complaint.
  \item \emph{Recall:} say the Japanese for not particularly, then the Japanese for a grumble, then say \emph{kujō} again
\end{itemize}

\begin{tcolorbox}[breakable,colback=teal!4,colframe=teal!35!black,title={Before you move on}]
What does \ja{くじょう} mean? (\textquotedblleft{}A formal complaint\textquotedblright{}.) Say it once more.
\end{tcolorbox}
\section[\texorpdfstring{\ja{ふへい} (fuhei)}{ふへい (fuhei)}]{\ja{ふへい} (fuhei) — discontent; grumbling}

\label{lesson:JA-C324-fuhei}



\begin{quote}
Before the new one: say the Japanese for a grumble, then the Japanese for a formal complaint.
\end{quote}

\subsection*{You'll want to know: \ja{ふへい}}
\textbf{\ja{ふへい}} — \emph{fuhei} — \textquotedblleft{}discontent; grumbling\textquotedblright{}.

\begin{quote}
\textbf{\ja{しごとが} \ja{おおいので、みんな} \ja{ふへいを} \ja{言}\ja{っています。}} — \emph{shigoto ga ōi node, minna fuhei o itte imasu} — \textquotedblleft{}There is a lot of work, so everyone is grumbling.\textquotedblright{}
\end{quote}

\begin{cousinweb}
\textbf{\ja{ふへい}} (\emph{fuhei}) is Sino-Japanese (kango): two Chinese characters read with their on'yomi, readings borrowed from Chinese. They are \emph{fu}, \textquotedblleft{}negative\textquotedblright{} (Middle Chinese \emph{pjut}) and \emph{hei}, \textquotedblleft{}even\textquotedblright{} (\emph{bjaeng}) (KANJIDIC2; Middle Chinese in Baxter's transcription).
\end{cousinweb}

\begin{grammarlens}[title={What you've built}]
One more word for giving reasons and opinions.
\end{grammarlens}

\subsection*{Your turn}
\begin{itemize}
\raggedright
  \item \emph{Say it:} \emph{fuhei}
  \item \emph{Say it:} \emph{fuhei}, once more
  \item \emph{Say it:} \emph{shigoto ga ōi node, minna fuhei o itte imasu} — There is a lot of work, so everyone is grumbling.
  \item \emph{Recall:} say the Japanese for a grumble, then the Japanese for a formal complaint, then say \emph{fuhei} again
\end{itemize}

\begin{tcolorbox}[breakable,colback=teal!4,colframe=teal!35!black,title={Before you move on}]
What does \ja{ふへい} mean? (\textquotedblleft{}Discontent; grumbling\textquotedblright{}.) Say it once more.
\end{tcolorbox}
\section[\texorpdfstring{\ja{ひはんする} (hihan suru)}{ひはんする (hihan suru)}]{\ja{ひはんする} (hihan suru) — to criticize}

\label{lesson:JA-C324-hihan-suru}



\begin{quote}
Before the new one: say the Japanese for a formal complaint, then the Japanese for discontent.
\end{quote}

\subsection*{You'll want to know: \ja{ひはんする}}
\textbf{\ja{ひはんする}} — \emph{hihan suru} — \textquotedblleft{}to criticize\textquotedblright{}.

\begin{quote}
\textbf{\ja{ねだんが} \ja{あがったので、みんな} \ja{かいしゃを} \ja{ひはんしました。}} — \emph{nedan ga agatta node, minna kaisha o hihan shimashita} — \textquotedblleft{}Prices went up, so everyone criticized the company.\textquotedblright{}
\end{quote}

Polite \textbf{\ja{ひはんします}}, \emph{hihanshimasu}; a \ja{する} verb. Past (polite): \textbf{\ja{しんぶんが} \ja{ひはんしました。}}, \emph{shinbun ga hihan shimashita}.

\begin{cousinweb}
\textbf{\ja{ひはんする}} (\emph{hihan suru}) is built from a Sino-Japanese noun and the native verb \emph{suru}, \textquotedblleft{}to do\textquotedblright{}. The noun, \emph{hihan}, is two Chinese characters read with their on'yomi, readings borrowed from Chinese: \emph{hi}, \textquotedblleft{}criticism\textquotedblright{} (Middle Chinese \emph{phej}) and \emph{han}, \textquotedblleft{}judgment\textquotedblright{} (\emph{phanH}) (KANJIDIC2; JMdict; Middle Chinese in Baxter's transcription).
\end{cousinweb}

\begin{grammarlens}[title={What you've built}]
One more word for giving reasons and opinions.
\end{grammarlens}

\subsection*{Your turn}
\begin{itemize}
\raggedright
  \item \emph{Say it:} \emph{hihan suru}
  \item \emph{Say it:} \emph{hihan suru}, once more
  \item \emph{Say it:} \emph{nedan ga agatta node, minna kaisha o hihan shimashita} — Prices went up, so everyone criticized the company.
  \item \emph{Recall:} say the Japanese for a formal complaint, then the Japanese for discontent, then say \emph{hihan suru} again
\end{itemize}

\begin{tcolorbox}[breakable,colback=teal!4,colframe=teal!35!black,title={Before you move on}]
What does \ja{ひはんする} mean? (\textquotedblleft{}To criticize\textquotedblright{}.) Say it once more.
\end{tcolorbox}
\section[\texorpdfstring{\ja{いいわけ} (iiwake)}{いいわけ (iiwake)}]{\ja{いいわけ} (iiwake) — an excuse}

\label{lesson:JA-C324-iiwake}



\begin{quote}
Before the new one: say the Japanese for discontent, then the Japanese for to criticize.
\end{quote}

\subsection*{You'll want to know: \ja{いいわけ}}
\textbf{\ja{いいわけ}} — \emph{iiwake} — \textquotedblleft{}an excuse\textquotedblright{}.

\begin{quote}
\textbf{\ja{おとうとは} \ja{おくれたので、いいわけを} \ja{しました。}} — \emph{otōto wa okureta node, iiwake o shimashita} — \textquotedblleft{}My younger brother was late, so he made an excuse.\textquotedblright{}
\end{quote}

\begin{cousinweb}
\textbf{\ja{いいわけ}} (\emph{iiwake}) is a native Japanese word (wago), inherited rather than borrowed from Chinese: it is built from \emph{ii}, \textquotedblleft{}saying (iu)\textquotedblright{} and \emph{wake}, \textquotedblleft{}reason\textquotedblright{}, and its kanji were attached later as kun'yomi (JMdict).
\end{cousinweb}

\begin{grammarlens}[title={What you've built}]
One more word for giving reasons and opinions.
\end{grammarlens}

\subsection*{Your turn}
\begin{itemize}
\raggedright
  \item \emph{Say it:} \emph{iiwake}
  \item \emph{Say it:} \emph{iiwake}, once more
  \item \emph{Say it:} \emph{otōto wa okureta node, iiwake o shimashita} — My younger brother was late, so he made an excuse.
  \item \emph{Recall:} say the Japanese for discontent, then the Japanese for to criticize, then say \emph{iiwake} again
\end{itemize}

\begin{tcolorbox}[breakable,colback=teal!4,colframe=teal!35!black,title={Before you move on}]
What does \ja{いいわけ} mean? (\textquotedblleft{}An excuse\textquotedblright{}.) Say it once more.
\end{tcolorbox}
